\documentclass[12pt]{article}
\usepackage[utf8]{inputenc}
\usepackage[T1]{fontenc}
\usepackage{hyperref}
\title{Wik-Iiyanh Language Resources for K-12 Curriculum Development}
\author{Compiled by Claude Sonnet 5.5 (claude-sonnet-5-5) for Andr\'as Kornai}
\date{}
\begin{document}
\maketitle

\section{Introduction}
Wik-Iiyanh (Wik Iyanh, Wik Iyanja, Wikianji; ISO 639-3 wij) is a language of central Cape York Peninsula, Queensland, southwest of Coen. Sources disagree on classification. Wikipedia treats Iiyanh as a dialect of Kugu-Muminh (Kugu Nganhcara), a Paman language, while Glottolog lists it within the Wik group. Joshua Project gives about 70--80 people, citing Ethnologue 2016. No current speaker count was found.

\section{Primary Educational Resources}
No primer, textbook, or school programme was found.

\section{Linguistic and Reference Materials}
Wikipedia entry: \url{https://en.wikipedia.org/wiki/Wik_Iyanh}. Related dialect cluster: \url{https://en.wikipedia.org/wiki/Kugu_Nganhcara_language}. Glottolog lists a reference by David F. Martin (1999): \url{https://glottolog.org/resource/reference/id/476984}. Its content was not checked.

\section{Religious and Community Resources}
Joshua Project profile: \url{https://joshuaproject.net/people_groups/15398}. No scripture was found.

\section{Digital and Audio Resources}
Global Recordings Network lists the language as verified but has no recordings: \url{https://globalrecordings.net/en/language/wij}.

\section{Curriculum Development Context}
Very little documentation was found. Any curriculum would need to rely on community knowledge and archival material still to be located, and possibly on materials from neighbouring Wik and Kugu languages after consultation with the community.
\end{document}